\documentclass{article}
\usepackage{tcee-repaso}
\tcetema[Tratados e instituciones de la UE]{3.B.36}{Tratados, orden jurídico e instituciones de la Unión Europea.}
\tcefecha{14/11/2026}
\begin{document}
\cabecerarepaso
\begin{columnas}

% ==========================================================================
\bloque{Esquema del cante}
\esquemacante{%
\apartadocante{Introducción: «del dicho al hecho hay mucho trecho»}{2,5}
\apartadocante{I. Los tratados: de París a Lisboa}{8}
\subcante{1-4. De 1945 a 1990}{3,5}
\subcante{5. De la Comunidad a la Unión (1990-2007)}{2}
\subcante{6. Lisboa, Brexit y ampliación}{2,5}
\apartadocante{II. El ordenamiento jurídico}{9}
\subcante{1. Bases: primacía, efecto directo, fuentes}{3,5}
\subcante{2. Competencias y principios (modelo)}{3,5}
\subcante{3. Aprobación de la legislación}{2}
\apartadocante{III. Las instituciones}{9,5}
\subcante{1-2. Comisión y Parlamento}{3}
\subcante{3. Consejo (poder de voto)}{3}
\subcante{4-8. Consejo Europeo, TJUE, TCu, BCE, otros}{3,5}
\apartadocante{Conclusión}{1}}
\emph{Gráficos del cante}: estructura de los tratados (\hyperref[rep:I]{I}), asignación de competencias y procedimiento legislativo (\hyperref[rep:II]{II}), poder de voto (\hyperref[rep:III]{III}); fórmulas: $e>\Delta/2$ y el índice de \autor{Shapley}-\autor{Shubik}.

% ==========================================================================
\bloque[2,5]{Introducción}
\begin{itemize}
\item \textbf{Enganche}: «Del dicho al hecho hay mucho trecho». \autor{Schuman} (9-5-1950): Europa se hará con realizaciones concretas que creen una solidaridad de hecho. Hoy: 27 Estados, euro en 21, Derecho que prima, 9 candidatos; Ucrania y Moldavia abren capítulos (junio de 2026). ¿Cómo convirtió la UE los dichos en hechos, y le sirven sus reglas con más de 30?
\item \textbf{Relevancia}: la integración más profunda del mundo; competencias exclusivas en comercio, moneda (euro) y competencia. España: miembro desde 1986 \vertema{4.A.27}, 10,9\,\% de la población, 61 de 720 escaños, $\approx$9\,\% del poder de voto.
\item \textbf{Contextualización}: \concepto{UE}: organización \emph{sui generis} de integración, personalidad jurídica desde Lisboa (2009). Comunidad de Derecho: \emph{ordenamiento propio}, \emph{instituciones propias}, \emph{atribución de competencias}. Presupuesto \vertema{3.B.37}; mercado interior \vertema{3.B.39}; euro \vertema{3.B.43}.
\item \textbf{Problemática}: ¿de 6 a 27? ¿Qué Derecho y qué competencias? ¿Qué nivel de gobierno para cada cosa? ¿Quién decide y con cuánto poder?
\item \textbf{Estructura}: \hyperref[rep:I]{I} historia; \hyperref[rep:II]{II} Derecho; \hyperref[rep:III]{III} instituciones.
\end{itemize}

% ==========================================================================
\bloque[8]{I. Los tratados\label{rep:I}}
\begin{ideaclave}
La integración avanzó a saltos, tras crisis y por la economía, alternando \emph{profundizar} (competencias, mayoría) y \emph{ampliar}; Lisboa ordenó el resultado y la ampliación exige nuevas reformas.
\end{ideaclave}

\sub{I.1-I.4. De 1945 a 1990\label{rep:I.1}}
\begin{itemize}
\item \emph{Antecedentes}: Plan Marshall (1948) $\to$ OECE (1948) y UEP (1950); Congreso de La Haya (1948): federalistas, intergubernamentalistas, funcionalistas (\autor{Monnet}); Consejo de Europa (1949).
\item \emph{París} (1951; vigor 1952-2002): CECA de los Seis; Alta Autoridad supranacional, Consejo, Asamblea, Tribunal. CED fracasa (1954).
\item \emph{Roma} (25-3-1957; vigor 1-1-1958): CEE (unión aduanera, completada el 1-7-1968; mercado común; políticas agrícola, comercial, transportes, competencia; FSE, BEI) y Euratom. Fusión de ejecutivos (1965; vigor 1967). Silla vacía (1965) $\to$ Compromiso de Luxemburgo (1966). 1.ª ampliación (1973): RU, IE, DK.
\item \emph{Crisis (1973-1985)}: Werner (1970), serpiente (1972), SME (1979) \vertema{3.B.43}; euroesclerosis. Pero: recursos propios (1970), Consejo Europeo (1974), FEDER (1975), Tribunal de Cuentas (1975/77), elecciones al PE (1979); Grecia (1981), España y Portugal (1986).
\item \emph{Relanzamiento}: Libro Blanco de \autor{Delors} (1985); \concepto{Acta Única} (1986; vigor 1987): mercado interior para 1992, mayoría cualificada, cooperación con el PE, nuevas competencias; Informe Delors (1989).
\end{itemize}

\sub{I.5. De la Comunidad a la Unión (1990-2007)\label{rep:I.5}}
\begin{itemize}
\item \concepto{Maastricht} (1992; vigor 1-11-1993): UE con 3 \emph{pilares} (CE supranacional; PESC y JAI intergubernamentales); UEM; ciudadanía; codecisión; subsidiariedad; Fondo de Cohesión; CdR. EEE (1994); AT, FI, SE (1995).
\item \emph{Ámsterdam} (1997; vigor 1999): Schengen al acervo; empleo; Alto Representante; \emph{cooperación reforzada}.
\item \emph{Niza} (2001; vigor 2003): triple mayoría; un comisario por Estado; ampliación de 2004 (10) y 2007 (BG, RO).
\item \emph{Constitución} (2004): rechazada por Francia y Países Bajos (2005).
\end{itemize}

\graficorepaso[\linewidth]{estructura-tratados}{3 pilares (Maastricht) $\to$ TUE + TFUE + Carta, con personalidad jurídica (Lisboa)}

\sub{I.6. Lisboa, Brexit y ampliación\label{rep:I.6}}
\begin{itemize}
\item \concepto{Lisboa} (13-12-2007; vigor 1-12-2009; Irlanda: no en 2008, sí en 2009): reforma, no texto nuevo. \concepto{TUE}, \concepto{TFUE}, \concepto{Carta} (mismo valor, art. 6) y Euratom. Aporta: personalidad jurídica (art. 47), fin de pilares, catálogo de competencias, PLO general, doble mayoría (2014), Consejo Europeo y BCE instituciones, presidente estable, AR/VP, iniciativa ciudadana, art. 50. Croacia (2013).
\item \emph{Brexit}: referéndum 23-6-2016 (51,9\,\%); notificación 29-3-2017; salida 31-1-2020; transición hasta 31-12-2020; Acuerdo de Comercio y Cooperación (2021); Marco de Windsor (2023).
\item \emph{Ampliación} (art. 49; criterios de Copenhague, 1993): 9 candidatos; Montenegro 21/33 capítulos (oct. 2026); Ucrania y Moldavia abren el bloque de fundamentales (junio de 2026) tras levantar Hungría el veto.
\item \emph{Reforma}: PE (22-11-2023, art. 48: más mayoría, Comisión menor); Comisión (6-10-2026): pasarelas (sanciones, fraude fiscal, energía), cooperación reforzada, salvaguardias de valores, sin reformar Tratados.
\end{itemize}

\graficorepaso[\linewidth]{ampliaciones}{De 6 a 28 y 27 Estados: ampliaciones y tratados, 1952-2026}

% ==========================================================================
\bloque[9]{II. El ordenamiento jurídico\label{rep:II}}
\begin{ideaclave}
Un ordenamiento propio que prima y se aplica directamente (lo dijo el TJUE, no los Tratados); la UE actúa solo donde se le atribuyen competencias y, fuera de las exclusivas, solo si lo hace mejor: la subsidiariedad es una regla económica.
\end{ideaclave}

\sub{II.1. Bases del ordenamiento\label{rep:II.1}}
\begin{itemize}
\item \emph{Rasgos}: autonomía (art. 344); \concepto{efecto directo} (\emph{Van Gend en Loos}, 1963); \concepto{primacía} (\emph{Costa/ENEL}, 1964; sobre constituciones, \emph{Internationale Handelsgesellschaft}, 1970; inaplicación, \emph{Simmenthal}, 1978); unidad (acervo); responsabilidad del Estado (\emph{Francovich}, 1991). Directivas: efecto vertical (\emph{Van Duyn}, 1974), no horizontal (\emph{Marshall}, 1986).
\item \emph{Conflictos}: BVerfG PSPP (2020) e infracción a Alemania (2021); Tribunal Constitucional polaco (2021).
\item \emph{Valores} (art. 2) y \emph{art. 7}: riesgo (4/5 del Consejo) y violación (unanimidad del Consejo Europeo); Polonia (2017-2024), Hungría (2018); condicionalidad (Reglamento 2020/2092; Hungría, 2022).
\item \emph{Fuentes}: originario (TUE, TFUE, Carta, Euratom) $>$ principios generales $>$ acuerdos internacionales (exclusivos o mixtos) $>$ derivado.
\item \emph{Derivado} (art. 288): \concepto{reglamento} (general, completo, directo: unificar); \concepto{directiva} (resultado, transposición: armonizar); \concepto{decisión} (destinatarios); recomendación y dictamen (no vinculantes). Legislativo según el procedimiento (art. 289). \emph{Delegados} (art. 290: elementos no esenciales; objeción del PE o del Consejo) y de \emph{ejecución} (art. 291; comitología, Reglamento 182/2011). \emph{Soft law}.
\end{itemize}

\graficorepaso[\linewidth]{actos-anuales}{55-150 actos del PLO al año, con picos al final de cada legislatura; $>$1.000 de ejecución}

\sub{II.2. Competencias y principios\label{rep:II.2}}
\begin{itemize}
\item \emph{Exclusivas} (art. 3): unión aduanera, competencia, política monetaria del euro, conservación de recursos marinos, comercial. \emph{Compartidas} (art. 4): mercado interior, cohesión, agricultura, medio ambiente, energía, transportes\ldots \emph{Coordinación} (art. 5), \emph{PESC} y \emph{apoyo} (art. 6: salud, industria, cultura, educación\ldots). Flexibilidad: art. 352; AETR (1971).
\item \emph{Principios} (art. 5 TUE): \concepto{atribución} (¿hay competencia?); \concepto{subsidiariedad} y \concepto{proporcionalidad} (¿debe ejercerla, cómo?); \emph{cooperación leal} (art. 4.3). Protocolo 2: tarjeta amarilla (1/3 de los votos de los Parlamentos nacionales) y naranja; 3 amarillas (2012, 2013, 2016).
\end{itemize}

\textbf{Asignación de competencias (\autor{Oates}, 1972; \autor{Alesina} \emph{et al.}, 2005)}
\begin{itemize}
\item \emph{Idea}: centralizar internaliza externalidades y aprovecha escala, pero uniformiza preferencias distintas \vertema{4.B.14}.
\item \emph{Supuestos}: 2 Estados, $\mathit{BMg}_i=a_i-G_i$, $\Delta=a_2-a_1$; coste $c$; externalidad $e$; la UE elige un $G$ uniforme.
\item \emph{Desarrollo}: óptimo $G_i^s=a_i-c+e$; descentralizado $G_i^D=a_i-c$, pérdida $e^2$; centralizado $G^C=\bar a-c+e$, pérdida $\Delta^2/4$.
\begin{formula}
\funcionclave{\text{centralizar}\iff e>\frac{\Delta}{2}}
\end{formula}
\item \emph{Implicaciones}: $e=0$: teorema de Oates; subsidiariedad = «dimensión o efectos»; exclusivas con $e$ alto y $\Delta$ bajo; ampliar sube $\Delta$.
\item \emph{Extensiones}: escala $F$ ($e^2+F>\Delta^2/4$); tamaño de la unión (\autor{Alesina} y \autor{Spolaore}); \autor{Tiebout}; segunda generación.
\item \emph{Evidencia}: la UE creció donde debía (comercio, mercado interior) y donde no (agricultura, regional), poco en defensa; Eurobarómetro.
\item \emph{Valoración}: criterio claro; la realidad responde a la historia y la política, y a reglas de voto concretas.
\end{itemize}

\pareja{oates-uniforme}{Uniformidad: exceso y defecto frente a $G_i^*$}{asignacion-competencias}{Frontera $e=\Delta/2$; con escala, más baja}

\sub{II.3. Aprobación de la legislación\label{rep:II.3}}
\begin{itemize}
\item \concepto{PLO} (arts. 289 y 294): Comisión propone (iniciativa casi exclusiva; PE, Consejo e iniciativa ciudadana: 1 millón, 7 Estados, pueden pedirla); 1.ª lectura PE y Consejo (sin plazo); 2.ª (3+1 meses): el PE aprueba, rechaza (mayoría de miembros) o enmienda; conciliación y 3.ª (6+2 semanas). Trílogos: casi todo en 1.ª lectura.
\item \emph{Especiales}: consulta (fiscalidad) y aprobación (MFP, art. 352, adhesiones). Presupuesto: art. 314 \vertema{3.B.37}.
\item \emph{Flexibilidad}: cooperación reforzada ($\geq$9 Estados; divorcio, patente unitaria, regímenes matrimoniales, Fiscalía Europea); PESCO (26); pasarelas (art. 48.7, nunca usada).
\end{itemize}

\graficorepaso[\linewidth]{procedimiento-legislativo}{1.ª lectura, 2.ª lectura, conciliación y 3.ª lectura}

\graficorepaso[\linewidth]{circulos-europa}{Geometría variable: UE, euro (21), Schengen (29), EEE, AELC, Consejo de Europa}

% ==========================================================================
\bloque[9,5]{III. Las instituciones\label{rep:III}}
\begin{ideaclave}
7 instituciones (art. 13 TUE) que equilibran ciudadanos (PE), Estados (Consejo, Consejo Europeo) e interés común (Comisión, TJUE, BCE); legislan el PE y el Consejo a propuesta de la Comisión.
\end{ideaclave}

\graficorepaso[\linewidth]{mapa-instituciones}{Quién elige a quién y el triángulo legislativo}

\sub{III.1-III.2. Comisión y Parlamento\label{rep:III.1}}
\begin{itemize}
\item \concepto{Comisión} (art. 17 TUE): 27 comisarios, 5 años; presidente propuesto por el Consejo Europeo (mayoría cualificada) y elegido por el PE (mayoría de miembros); Colegio aprobado por el PE. \autor{Von der Leyen} II desde 1-12-2024; 6 vicepresidentes ejecutivos (\autor{Ribera}; AR \autor{Kallas}). Funciones: iniciativa, ejecución (presupuesto, actos delegados y de ejecución), guardiana (infracciones, arts. 258-260; competencia), representación exterior.
\item \concepto{Parlamento} (art. 14 TUE): máximo 750 + presidente; 720 (2024-2029), proporcionalidad decreciente (6-96; España 61); grupos: PPE 188, S\&D 136, PfE 84, ECR 78, Renew 77, Verdes 53, Izquierda 46, ESN 25, NI 33 (julio de 2024). Presidenta \autor{Metsola}. Funciones: legislativa, presupuestaria (descargo), control (censura: 2/3 de votos y mayoría de miembros).
\end{itemize}

\graficorepaso[\linewidth]{escanos-estados}{Proporcionalidad decreciente: 870.000 habitantes por escaño en Alemania, 96.000 en Malta}

\graficorepaso[0.85\linewidth]{grupos-pe}{Ninguna mayoría: PPE + S\&D + Renew = 401}

\sub{III.3. El Consejo de la UE\label{rep:III.3}}
\begin{itemize}
\item Ministros; 10 formaciones; presidencia semestral en tríos (Chipre y, desde julio de 2026, Irlanda; luego Lituania y Grecia); Coreper de embajadores y de adjuntos.
\item \emph{Mayoría cualificada}: $\geq$55\,\% de Estados (15) y $\geq$65\,\% de la población (293,9 de 452,2 millones en 2026); 72\,\% si no propone la Comisión; minoría de bloqueo $\geq$4. Unanimidad: PESC, fiscalidad, MFP, adhesiones.
\end{itemize}

\textbf{Poder de voto (\autor{Shapley} y \autor{Shubik}, 1954; \autor{Banzhaf}, 1965)}
\begin{itemize}
\item \emph{Idea}: el peso no es el poder; importa ser decisivo \vertema{3.A.14} \vertema{4.B.2}.
\item \emph{Supuestos}: 27 votantes; gana $T$ si $|T|\geq15$ y $\geq$65\,\% o $|T|\geq24$; a priori.
\item \emph{Desarrollo}: pivote si $v(T\cup\{i\})-v(T)=1$.
\begin{formula}
\funcionclave{\phi_i(v)=\sum_{T\subseteq N\setminus\{i\}}\frac{|T|!\,(n-|T|-1)!}{n!}\big[v(T\cup\{i\})-v(T)\big]}
\end{formula}
Banzhaf: proporción de coaliciones en que es decisivo. Alemania: población 18,5, SS 17,5, Bz 12,0; España 10,9, 9,4, 7,8; Malta 0,13, 0,87, 1,73.
\item \emph{Implicaciones}: 2 legitimidades; pequeños sobrerrepresentados; probabilidad de aprobar 13,2\,\% (Niza $\approx$2\,\%); España perdió frente a Alemania.
\item \emph{Extensiones}: raíz cuadrada (\autor{Penrose}, 1946; Jagellón); poder a posteriori; ampliación.
\item \emph{Evidencia}: $\approx$80\,\% por mayoría cualificada, pero casi todo sin votos en contra; Niza y la Constitución se negociaron con estos índices.
\item \emph{Valoración}: útiles para diseñar reglas; ignoran preferencias, agenda y PE.
\end{itemize}

\graficorepaso[0.85\linewidth]{poder-voto}{Por encima de la bisectriz, más poder que población}

\sub{III.4-III.8. Las demás\label{rep:III.4}}
\begin{itemize}
\item \concepto{Consejo Europeo} (art. 15): impulso político, no legisla; consenso; presidente 2,5 años renovable (\autor{Costa}, 1-12-2024 a 31-5-2027); nombramientos.
\item \concepto{TJUE} (art. 19; Luxemburgo): Tribunal de Justicia (27 jueces, 11 abogados generales; Gran Sala de 15) y Tribunal General (54; prejudiciales en 6 ámbitos desde 2024); cuestión prejudicial (267), incumplimiento (258-260), anulación (263), omisión (265). Tribunal de la Función Pública suprimido (2016).
\item \concepto{Tribunal de Cuentas} (art. 285): auditor externo, uno por Estado, 6 años; informe anual sobre el año anterior y declaración de fiabilidad.
\item \concepto{BCE} (arts. 282-284; Fráncfort): estabilidad de precios (2\,\% simétrico); Comité Ejecutivo (\autor{Lagarde}; \autor{Vujčić} VP desde 1-6-2026), Consejo de Gobierno (21 gobernadores), Consejo General \vertema{3.B.43}.
\item \emph{Otros}: CESE y CdR (329 cada uno; 21 de España); BEI (\autor{Calviño}) \vertema{3.B.30}; Defensor del Pueblo; Fiscalía Europea (25 Estados con Hungría, 2026; \autor{Ritter} desde 1-11-2026); SEAE; $>$30 agencias; Eurogrupo (\autor{Pierrakakis}, 12-12-2025); MEDE (intergubernamental, 500.000 M€) \vertema{3.B.45}.
\end{itemize}

% ==========================================================================
\bloque[1]{Conclusión}
\begin{itemize}
\item \textbf{Recapitulación}: realizaciones concretas: de 6 a 27 y del carbón a la moneda; Derecho que prima; subsidiariedad como regla económica; instituciones que equilibran 3 legitimidades.
\item \textbf{Relevancia}: el marco de casi toda la política económica española \vertema{4.A.30}.
\item \textbf{Extensiones}: \vertema{3.B.10} \vertema{3.B.37} \vertema{3.B.39} \vertema{3.B.41} \vertema{3.B.43} \vertema{3.B.45} \vertema{4.B.14}.
\item \textbf{Opinión}: ampliar con pasarelas, cooperación reforzada y salvaguardias; la reforma completa, después.
\item \textbf{Idea final}: del dicho al hecho hay mucho trecho: cada tratado convirtió un dicho en un hecho jurídico; para pasar de 30, otra vez hay que convertir en reglas lo prometido.
\end{itemize}

% ==========================================================================
\bloque{Autores y fechas}
\begin{itemize}
\item 1946 \autor{Penrose}; 1948 La Haya; 1950 \autor{Schuman}; 1951 París; 1954 \autor{Shapley}-\autor{Shubik}; 1957 Roma; 1958 \autor{Haas}; 1963 \emph{Van Gend en Loos}; 1964 \emph{Costa/ENEL}; 1965 \autor{Banzhaf} y fusión; 1966 Luxemburgo; 1972 \autor{Oates}; 1978 \emph{Simmenthal}; 1986 Acta Única; 1991 \emph{Francovich}; 1992 Maastricht; 1997 Ámsterdam; 1998 \autor{Moravcsik}; 2001 Niza; 2004 Constitución; 2005 \autor{Alesina}, \autor{Angeloni} y \autor{Schuknecht}; 2007 Lisboa; 2014 doble mayoría; 2020 Brexit; 2023 propuesta del PE; 2026 Ucrania y Moldavia abren capítulos.
\end{itemize}

\bloque{Datos clave}
\begin{itemize}
\item 27 Estados, 452,2 millones de habitantes (2026); euro en 21 (Bulgaria, 2026); Schengen 29.
\item PE: 720 escaños; España 61. Consejo: umbral de 293,9 millones.
\item Actos 2025: 55 del PLO, 166 delegados, 1.203 de ejecución.
\item Candidatos: 9 + Kosovo; Montenegro 21 de 33 capítulos (oct. 2026).
\end{itemize}

\bloque{Preguntas del tribunal}
\begin{enumerate}
\item \textbf{¿Qué distingue a la UE de otras organizaciones internacionales?} Primacía, efecto directo, mayoría, PE elegido, TJUE: supranacionalidad.
\item \textbf{¿Consejo Europeo, Consejo de la UE y Consejo de Europa?} Jefes de Estado o de Gobierno; ministros; organización de 46 Estados (derechos humanos).
\item \textbf{¿Qué es una minoría de bloqueo?} Al menos 4 Estados con más del 35\,\% de la población.
\item \textbf{¿Por qué no basta la mayoría simple?} 2 legitimidades (Estados y ciudadanos); protección de pequeños y grandes.
\item \textbf{Subsidiariedad y proporcionalidad}: si debe actuar la UE y cómo; tarjetas amarilla y naranja.
\item \textbf{¿Reglamento o directiva?} Unificar o armonizar; efecto directo de las directivas solo vertical.
\item \textbf{¿Qué papel tiene el PE en los acuerdos con terceros?} Aprobación (art. 218.6); en los comerciales, imprescindible \vertema{3.B.41}.
\item \textbf{¿Qué aportó Maastricht y cómo era su estructura?} 3 pilares, UEM, ciudadanía, codecisión.
\item \textbf{¿Quién elabora, aprueba, ejecuta y controla el presupuesto?} Comisión; Consejo y PE; Comisión; Tribunal de Cuentas y PE (descargo) \vertema{3.B.37}.
\item \textbf{¿Gigante económico, enano político?} PESC por unanimidad, defensa nacional; avances: PESCO, Fondo Europeo de Defensa, AR/VP.
\item \textbf{¿Puede un tribunal nacional desobedecer al TJUE?} No (art. 344, primacía); casos PSPP (2020) y Polonia (2021).
\item \textbf{¿Hay que reformar los Tratados para ampliar?} No necesariamente: pasarelas, cooperación reforzada, salvaguardias (Comisión, 2026).
\end{enumerate}

\begin{trampa}
\begin{itemize}
\item París: firmado en \emph{1951}, en vigor en 1952, expiró en \emph{2002}. Roma: 25-3-1957, vigor 1-1-1958.
\item Unión aduanera completada el \emph{1-7-1968}; España y Portugal, \emph{1-1-1986}; Croacia, \emph{2013}.
\item Lisboa: firmado el 13-12-2007, en vigor el \emph{1-12-2009}; doble mayoría desde el \emph{1-11-2014}.
\item La primacía \emph{no} está en los Tratados (Declaración 17); la norma nacional contraria se \emph{inaplica}, no es nula.
\item Cooperación reforzada: al menos \emph{9} Estados (Niza: 8).
\item Tribunal de Cuentas: institución desde \emph{Maastricht}; Consejo Europeo y BCE, desde \emph{Lisboa}.
\item Gran Sala del TJ: \emph{15} jueces; Tribunal General: \emph{54}.
\item Censura: \emph{2/3} de los votos emitidos y mayoría de los miembros del PE.
\end{itemize}
\end{trampa}

\end{columnas}
\end{document}
